% Figure: schematic of the Mars Climate Orbiter trajectory error
% accumulation across nine months of cruise.
\begin{figure}[htbp]
\centering
\begin{tikzpicture}[
  axis/.style={draw=engblack, -{Stealth}, thick},
  planned/.style={draw=engblue, thick, dashed},
  actual/.style={draw=engred, thick},
  label/.style={font=\small},
]
% Axes
\draw[axis] (0,0) -- (10,0) node[anchor=west, label] {time (months)};
\draw[axis] (0,0) -- (0,5)  node[anchor=south, label] {periapsis altitude (km)};

% Tick marks on x-axis
\foreach \x/\m in {0/0, 1.5/2, 3/4, 4.5/6, 6/8, 7.5/10} {
  \draw (\x,-0.1) -- (\x,0.1);
  \node[font=\scriptsize, below] at (\x,-0.1) {\m};
}

% Tick marks on y-axis
\foreach \y/\h in {0/0, 1/50, 2/100, 3/150, 4/200, 5/250} {
  \draw (-0.1,\y) -- (0.1,\y);
  \node[font=\scriptsize, left] at (-0.1,\y) {\h};
}

% Planned trajectory (constant at 226 km)
\draw[planned] (0, 4.52) -- (8.0, 4.52);
\node[engblue, font=\small, anchor=west] at (8.0, 4.52)
  {\textbf{planned $\approx 226$ km}};

% Actual trajectory (linear decline from same start to 57 km at 9 months)
\draw[actual] (0, 4.52) -- (8.0, 1.14);
\node[engred, font=\small, anchor=west] at (8.0, 1.14)
  {\textbf{actual $\approx 57$ km}};

% Atmosphere band
\fill[engred, opacity=0.15] (0,0) rectangle (8.0, 2.0);
\node[engred, font=\small\itshape] at (4,0.5) {Mars atmosphere (loss)};

% Annotation
\node[font=\footnotesize, anchor=west, align=left] at (0.3, 3.0)
  {Per-event bias: $\sim 4.45\times$ overstatement\\
   in small-forces impulse, accumulating\\
   through angular-momentum-desaturation\\
   manoeuvres across cruise.};

\end{tikzpicture}
\caption[MCO trajectory error accumulation]{Schematic of the Mars
Climate Orbiter periapsis altitude over the nine-month cruise. The
planned closest-approach altitude was approximately $226 \,\si{\kilo
\meter}$. The accumulated bias from the small-forces unit
mismatch produced an actual closest approach of approximately
$57 \,\si{\kilo\meter}$, below the survivability threshold for
atmospheric loading. The error grew incrementally with each
small-impulse event over the cruise and was not large enough at any
single moment to trigger a navigation alarm. Reconstructed from
\cite{acc:nasa-mco-mib}.}
\label{fig:vol01:ch01:mco-trajectory}
\end{figure}
